\documentclass[11pt]{letter}

\usepackage[margin=0.75in]{geometry}
\usepackage{setspace}
\usepackage{hyperref}
\usepackage{xcolor}
\usepackage{enumitem}
\usepackage{parskip}

\hypersetup{
    colorlinks=true,
    urlcolor=black,
    linkcolor=black
}

\setlength{\parindent}{0pt}
\setstretch{1.05}
\date{}
\begin{document}

\begin{letter}{}

\opening{Dear Hiring Committee,}

I saw the job post on \textbf{Amazon Jobs} and thought I'd be a good fit for it.

I am a Master's student in Computer Science at the University of Cincinnati, graduating December 2026, looking for a role as a \textbf{Software Development Engineer} at \textbf{Amazon EBS}. With hands on experience owning the design, implementation, and operation of large-scale distributed AWS infrastructure, I am excited about the opportunity to contribute to \textbf{Amazon EBS} and support the EBS Stats team delivering fleet telemetry and performance insights at petabyte scale.

\textbf{Here are reasons that make me a great fit for the role:}

\begin{itemize}[leftmargin=*, itemsep=0.45em]
    \item \textbf{Full software development lifecycle ownership:} At Siemens Industry Software, I owned design and operation of Teamcenter X infrastructure workflows end to end, automating with Python (Boto3) and cutting deployment time 40\% across AWS environments, then executed 44 production environment upgrades in 4 months, diagnosing live infrastructure issues across Linux, EC2, and databases.

    \item \textbf{Distributed systems design and reliability:} I designed an Ansible based Blue-Green deployment framework that reduced production upgrade downtime by 50\% with reliable rollback, and migrated 40+ production environments from process-based deployments to Kubernetes using ArgoCD, Helm, GitLab CI/CD, and HashiCorp Vault, building the kind of fault-tolerant operational discipline EBS's storage fleet demands.

    \item \textbf{High-throughput data pipeline experience:} I built a distributed data pipeline with PySpark and Delta Lake processing a 2.3M-row hourly telemetry dataset across 263 zones, and engineered a 4-stage Python ETL pipeline processing 71K+ events with a PostgreSQL analytics schema and materialized views for low-latency queries, directly relevant to ingesting and querying fleet-scale telemetry.

    \item \textbf{Code quality, mentorship, and code review culture:} As Quality Coach and SonarQube Specialist for the Agile Release Train, I increased automated code coverage from 35\% to 78\% by enforcing coding standards and CI/CD quality gates, and mentored 5 software engineers through structured onboarding, reflecting the collaborative, kind-code-review culture EBS values.
\end{itemize}

I am available for full-time work after my graduation in December 2026. As a Master's student at the University of Cincinnati and someone who values professionalism and growth, I am enthusiastic about contributing to \textbf{Amazon EBS} and helping deliver reliable, low-latency insights from one of the largest distributed storage systems in the world.

Please find my CV attached below. I look forward to hearing from you. I would be excited to help EBS Stats turn petabyte-scale fleet telemetry into insights that make AWS storage better for millions of customers.

{\setlength{\parskip}{0pt}
Regards,\\
Tanmay Bhuskute,\\
tanmay.v.bhuskute@gmail.com $\mid$ +1 (732) 858-3027,\\
Seattle, WA.
}

\end{letter}

\end{document}
